\subsection{Neutral Generators and Character-Vanishing Conditions}
  \label{subsec:neutral-generators}

  The admissibility hierarchy observed in the previous sections suggests that
  not all generators contribute equally to dynamically sustainable
  configurations.
  A useful structural notion capturing this behaviour is that of
  \emph{neutral generators}.

  Let $\rho$ be an irreducible representation of a finite group $G$, and let
  $\chi_\rho$ denote its character.
  A generator $s\in G$ will be called \emph{neutral} in the sector $\rho$ if
  \begin{equation}
    \chi_\rho(s)=0 .
  \end{equation}

  This condition expresses a cancellation property at the level of the
  representation trace.
  Although the operator $\rho(s)$ itself need not vanish, the vanishing of its
  trace implies that its eigenvalues occur in balanced pairs whose sum is zero.
  Such elements therefore play a distinguished role in organising the geometry
  of admissible generating sets.

  Character-vanishing conditions are not exceptional.
  In irreducible representations of dimension greater than one, many conjugacy
  classes naturally satisfy $\chi_\rho(s)=0$.
  The condition therefore provides a representation-theoretic criterion for
  selecting generators that remain compatible with admissible spectral
  configurations.

  In concrete realisations of binary groups inside $SU(2)$, neutral generators
  correspond to elements whose representation matrices are traceless.
  These elements act as half-turn rotations in the spinorial representation.
  However, the definition of neutrality does not depend on this specific
  realisation and can be formulated entirely in terms of representation
  characters.

  Neutral generators impose strong constraints on the mutual relations between
  elements of the generating set.
  When several generators satisfy the character-vanishing condition in the same
  representation sector, their pairwise products determine a rigid pattern of
  character values.

  As we show in the following section, these relations can be encoded in a Gram
  matrix whose entries are determined by representation characters.
  This observation provides a bridge between spectral admissibility and the
  geometric structure of generating sets.
